\chapter{Bibliografia}
\label{ap:bibliografia}

Todas as referências da apostila, separadas por tipo. Cada registro diz como
obter a obra: \emph{à venda na editora}, \emph{acesso aberto},
\emph{acesso pago} (só o DOI) ou \emph{esgotado}, mais a edição brasileira e
a leitura gratuita quando existem. Onde achar cada uma, em tabela, está no
apêndice~\ref{ap:onde-achar}.

\printbibliography[keyword=livro, heading=subbibliography, title={Livros}]
\printbibliography[keyword=paper, heading=subbibliography, title={Papers e teses}]
\printbibliography[keyword=artigo, heading=subbibliography, title={Artigos e blog posts}]
\printbibliography[keyword=video, heading=subbibliography, title={Vídeos}]
\printbibliography[keyword=doc, heading=subbibliography, title={Documentação oficial}]
\printbibliography[keyword=ferramenta, heading=subbibliography, title={Ferramentas}]
\printbibliography[keyword=curso, heading=subbibliography, title={Cursos gratuitos}]
\printbibliography[keyword=comunidade, heading=subbibliography, title={Comunidades}]
